%==============================================================================
%  CHAPITRE 1 - Méthodologie : Scrum + DevOps
%==============================================================================
\chapter{Méthodologie : Scrum + DevOps}
\label{chap:methodologie}

\section{Objet du document}

Ce document constitue la \textbf{roadmap de développement} du projet Zümm.
Il traduit le cahier des charges (v1.0, 13 juillet 2026) en un plan
d'exécution opérationnel : un \emph{product backlog} priorisé, une
planification en \textbf{18 sprints} de deux semaines avec plans d'exécution
détaillés, un registre des risques coté, une chaîne d'intégration et de
déploiement continus (CI/CD), un plan de releases sémantiques et un
dispositif de monitoring.

\begin{encadre}
\textbf{Chiffres clés.} 19 epics --- 80 user stories --- 582 story points ---
1 sprint de cadrage + 18 sprints de 14 jours --- 4 releases
(v0.1.0 $\rightarrow$ v1.0.0) --- période : 14 juillet 2026 $\rightarrow$
26 avril 2027.
\end{encadre}

\section{Cadre Scrum}

\subsection{Rôles}

\begin{center}
\begin{tabular}{lp{10.5cm}}
\toprule
\textbf{Rôle} & \textbf{Responsabilité} \\
\midrule
Product Owner & Priorise le backlog, valide les critères d'acceptation,
                accepte les livrables en sprint review. \\
Scrum Master  & Garant du cadre, lève les obstacles, anime les cérémonies. \\
Équipe de développement & Conçoit, développe, teste et livre l'incrément
                de chaque sprint. \\
\bottomrule
\end{tabular}
\end{center}

\subsection{Cérémonies}

Chaque sprint suit le même rythme de cérémonies :

\begin{center}
\begin{tabular}{llp{6.6cm}}
\toprule
\textbf{Cérémonie} & \textbf{Cadence / durée} & \textbf{Objectif} \\
\midrule
Sprint Planning       & Jour 1, 09h00--13h00 (4 h) & Sélection des stories et
                        engagement de l'équipe. \\
Daily Scrum           & Tous les jours, 09h15 (15 min) & Synchronisation,
                        détection des blocages. \\
Sprint Review         & Dernier jour, 14h00--16h00 (2 h) & Démonstration de
                        l'incrément au Product Owner. \\
Sprint Retrospective  & Dernier jour, 16h00--17h00 (1 h) & Amélioration
                        continue du processus. \\
\bottomrule
\end{tabular}
\end{center}

\subsection{Definition of Ready (DoR)}

Une story ne peut entrer dans un sprint que si :
\begin{enumerate}
  \item elle a un titre clair et une description ;
  \item ses critères d'acceptation sont définis ;
  \item ses dépendances sont identifiées ;
  \item elle est estimée en points ;
  \item le Product Owner a validé sa priorité.
\end{enumerate}

\subsection{Definition of Done (DoD)}

Une story n'est terminée que si :
\begin{enumerate}
  \item le code est développé et revu (\emph{peer review}) ;
  \item les tests unitaires passent (couverture $\geq 70\,\%$) ;
  \item les tests d'intégration passent ;
  \item la documentation est mise à jour ;
  \item l'incrément est déployé sur l'environnement de staging ;
  \item le Product Owner a validé la story.
\end{enumerate}

\section{Principes DevOps}

La démarche DevOps complète Scrum sur l'axe technique :

\begin{itemize}
  \item \textbf{Intégration continue} : chaque \emph{push} déclenche build,
        tests et analyse qualité (chapitre~\ref{chap:devops}) ;
  \item \textbf{Déploiement continu} : livraison automatique en staging,
        déploiement en production sur validation manuelle (\emph{approval gate}) ;
  \item \textbf{Infrastructure as Code} : Docker, docker-compose,
        manifestes Kubernetes optionnels ;
  \item \textbf{Observabilité} : logs centralisés, métriques Prometheus,
        dashboards Grafana, alerting (chapitre~\ref{chap:monitoring}) ;
  \item \textbf{Boucle de rétroaction} : les métriques de production
        alimentent la priorisation du backlog.
\end{itemize}

\begin{encadre}
\textbf{Note de cohérence.} Le plan initial couvrait 8 sprints pour
304 points. Le périmètre s'est étendu au fil des revues --- durcissement de
la sécurité, PWA livrable, conformité miel, puis solde des dettes --- pour
atteindre \textbf{18 sprints de construction et 582 points}. La somme des
\texttt{story\_points} des dix-neuf sprints de \texttt{sprints.json} égale
le total du \emph{product backlog}, et les trois représentations --- product
backlog, fichiers opérationnels et présent document --- sont alignées.

\smallskip
\textbf{Réserve sur la vélocité annoncée.} Les 582 points comptent les 37
points du sprint 8, dont \textbf{21 relèvent de livrables documentaires}
(US-039 diagrammes UML, US-040 rapport et soutenance) qui ne sont pas
produits : la vélocité \emph{applicative} réellement livrée est donc de
\textbf{561 points}. Le compter autrement serait surévaluer le livrable.
\end{encadre}